
\begin{ex} % Câu 9
	Rút gọn các biểu thức sau: $B=\left(\frac{x-6}{x+3\sqrt{x}}+\frac{1}{\sqrt{x}+3}\right)\left(1-\frac{1}{\sqrt{x}-2}\right)$ với $x>0; x \ne 4$.
	\loigiai{
		Điều kiện: $x>0; x \ne 4$.
		Ta có $x+3\sqrt{x} = \sqrt{x}(\sqrt{x}+3)$.
		Xét phần trong ngoặc thứ nhất:
		\begin{align*}
			\frac{x-6}{x+3\sqrt{x}}+\frac{1}{\sqrt{x}+3} &= \frac{x-6}{\sqrt{x}(\sqrt{x}+3)}+\frac{\sqrt{x}}{\sqrt{x}(\sqrt{x}+3)} \\
			&= \frac{x-6+\sqrt{x}}{\sqrt{x}(\sqrt{x}+3)} \\
			&= \frac{(\sqrt{x}+3)(\sqrt{x}-2)}{\sqrt{x}(\sqrt{x}+3)} \\
			&= \frac{\sqrt{x}-2}{\sqrt{x}}
		\end{align*}
		Xét phần trong ngoặc thứ hai:
		\begin{align*}
			1-\frac{1}{\sqrt{x}-2} &= \frac{\sqrt{x}-2-1}{\sqrt{x}-2} \\
			&= \frac{\sqrt{x}-3}{\sqrt{x}-2}
		\end{align*}
		Vậy biểu thức $B$ trở thành:
		\begin{align*}
			B &= \left(\frac{\sqrt{x}-2}{\sqrt{x}}\right) \cdot \left(\frac{\sqrt{x}-3}{\sqrt{x}-2}\right)
		\end{align*}
		Vì $x \ne 4$ nên $\sqrt{x}-2 \ne 0$. Ta có thể rút gọn:
		\begin{align*}
			B &= \frac{\sqrt{x}-3}{\sqrt{x}}
		\end{align*}
		Vậy, với $x>0; x \ne 4$ thì $B = \frac{\sqrt{x}-3}{\sqrt{x}}$.
	}
\end{ex}

\begin{ex} % Câu 10
	Giải hệ phương trình: $ \begin{cases} 2x-y=-3 \\ x+3y=2 \end{cases} $
	\loigiai{
		Ta có hệ phương trình:
		$ \begin{cases} 2x-y=-3 \quad &(1) \\ x+3y=2 \quad &(2) \end{cases} $
		Từ phương trình (1), suy ra $y = 2x+3$.
		Thế $y = 2x+3$ vào phương trình (2):
		\begin{align*}
			x + 3(2x+3) &= 2 \\
			x + 6x + 9 &= 2 \\
			7x &= 2 - 9 \\
			7x &= -7 \\
			x &= -1
		\end{align*}
		Thay $x=-1$ vào $y = 2x+3$:
		\begin{align*}
			y &= 2(-1) + 3 \\
			y &= -2 + 3 \\
			y &= 1
		\end{align*}
		Vậy, nghiệm của hệ phương trình là $(x;y) = (-1;1)$.
	}
\end{ex}

\begin{ex} % Câu 11
	Gọi $x_1, x_2$ là hai nghiệm của phương trình $x^2+4x-6=0$. Không giải phương trình, tính giá trị của biểu thức sau: $P=\frac{x_1+2}{x_2}+\frac{x_2+2}{x_1}$.
	\loigiai{
		Phương trình đã cho là $x^2+4x-6=0$.
		Theo hệ thức Vi-ét, ta có:
		$ \begin{cases} x_1 + x_2 = -4 \\ x_1 x_2 = -6 \end{cases} $
		Ta biến đổi biểu thức P:
		\begin{align*}
			P &= \frac{x_1+2}{x_2}+\frac{x_2+2}{x_1} \\
			&= \frac{x_1(x_1+2) + x_2(x_2+2)}{x_1x_2} \\
			&= \frac{x_1^2+2x_1 + x_2^2+2x_2}{x_1x_2} \\
			&= \frac{(x_1^2+x_2^2) + 2(x_1+x_2)}{x_1x_2} \\
			&= \frac{((x_1+x_2)^2 - 2x_1x_2) + 2(x_1+x_2)}{x_1x_2}
		\end{align*}
		Thay giá trị từ hệ thức Vi-ét vào biểu thức P:
		\begin{align*}
			P &= \frac{((-4)^2 - 2(-6)) + 2(-4)}{-6} \\
			&= \frac{(16 + 12) - 8}{-6} \\
			&= \frac{28 - 8}{-6} \\
			&= \frac{20}{-6} \\
			&= -\frac{10}{3}
		\end{align*}
		Vậy, giá trị của biểu thức $P = -\frac{10}{3}$.
	}
\end{ex}
\begin{ex} % Câu 12
	Hưởng ứng phong trào trồng cây xanh vì môi trường xanh sạch đẹp, một chi đoàn thanh niên dự định trồng 400 cây trong một thời gian quy định. Mỗi ngày chi đoàn đã trồng vượt mức kế hoạch 10 cây. Do vậy, chi đoàn đã hoàn thành công việc sớm hơn thời gian quy định 2 ngày. Hỏi mỗi ngày chi đoàn đã trồng được bao nhiêu cây?
	\loigiai{
		Gọi $x$ là số cây dự định trồng mỗi ngày theo kế hoạch. (Điều kiện: $x > 0$, $x$ là số nguyên)
		Thời gian dự định hoàn thành công việc là: $\frac{400}{x}$ (ngày).
		
		Thực tế, mỗi ngày chi đoàn đã trồng được: $x+10$ (cây).
		Thời gian thực tế hoàn thành công việc là: $\frac{400}{x+10}$ (ngày).
		
		Theo đề bài, chi đoàn đã hoàn thành công việc sớm hơn thời gian quy định 2 ngày, nên ta có phương trình:
		$\frac{400}{x} - \frac{400}{x+10} = 2$
		
		Nhân cả hai vế của phương trình với $x(x+10)$ để khử mẫu (với $x \ne 0, x \ne -10$):
		\begin{align*}
			400(x+10) - 400x &= 2x(x+10) \\
			400x + 4000 - 400x &= 2x^2 + 20x \\
			4000 &= 2x^2 + 20x \\
			2x^2 + 20x - 4000 &= 0 \\
			x^2 + 10x - 2000 &= 0
		\end{align*}
		Giải phương trình bậc hai $x^2 + 10x - 2000 = 0$:
		$\Delta = b^2 - 4ac = 10^2 - 4(1)(-2000) = 100 + 8000 = 8100$
		$\sqrt{\Delta} = \sqrt{8100} = 90$\\
		Phương trình có hai nghiệm phân biệt:\\
		$x_1 = \frac{-10 + 90}{2 \cdot 1} = \frac{80}{2} = 40$ (Thỏa mãn điều kiện $x>0$)\\
		$x_2 = \frac{-10 - 90}{2 \cdot 1} = \frac{-100}{2} = -50$ (Không thỏa mãn điều kiện $x>0$)\\
		
		Vậy, số cây dự định trồng mỗi ngày theo kế hoạch là 40 cây.
		Số cây thực tế mỗi ngày chi đoàn đã trồng được là: $x+10 = 40+10 = 50$ (cây).
		
		Trả lời: Mỗi ngày chi đoàn đã trồng được 50 cây.
	}
\end{ex}
\begin{ex} % Câu 13
	Bạn An rút ngẫu nhiên một tấm thẻ từ trong hộp chứa 40 tấm thẻ ghi các số $5; 6; 7; \dots; 44$; hai thẻ khác nhau thì ghi hai số khác nhau. Tính xác suất của các biến cố A: “Rút được tấm thẻ ghi số có 2 chữ số giống nhau”.
	\loigiai{
		Tổng số tấm thẻ trong hộp là 40. Vậy, số phần tử của không gian mẫu là $n(\Omega) = 40$.
		Các số được ghi trên thẻ là từ 5 đến 44.
		
		Biến cố A: “Rút được tấm thẻ ghi số có 2 chữ số giống nhau”.
		Ta liệt kê các số có hai chữ số giống nhau trong khoảng từ 5 đến 44:
		\begin{itemize}
			\item Số 11
			\item Số 22
			\item Số 33
			\item Số 44
		\end{itemize}
		Có 4 số thỏa mãn điều kiện của biến cố A.
		Vậy, số phần tử của biến cố A là $n(A) = 4$.
		
		Xác suất của biến cố A là:
		$P(A) = \frac{n(A)}{n(\Omega)} = \frac{4}{40} = \frac{1}{10}$.
	}
\end{ex}
\begin{ex} % Câu 14
	Cho đường tròn $(O)$ đường kính $AB$ và điểm $M$ bất kì thuộc đường tròn sao cho $MA < MB$ ($M \ne A$). Kẻ tiếp tuyến tại $A$ của đường tròn, tiếp tuyến này cắt tia $BM$ ở $N$. Tiếp tuyến của đường tròn tại $M$ cắt $AN$ ở $D$.
	\begin{enumerate}[a)]
		\item Chứng minh bốn điểm $A, D, M, O$ cùng thuộc một đường tròn.
		\item Chứng minh $OD$ song song $BM$. Qua $O$ kẻ đường thẳng vuông góc với $AB$ và cắt đường thẳng $BM$ tại $I$. Gọi giao điểm của $AI$ và $BD$ là $G$. Chứng minh ba điểm $N, G, O$ thẳng hàng.
	\end{enumerate}
	\loigiai{}
\end{ex}

\begin{ex} % Câu 15
	\begin{enumerate}[a)]
		\item Một xe tải chở hàng có thùng hàng phía sau dạng hình hộp chữ nhật với chiều rộng $2,4\text{m}$ và chiều cao tính từ mặt đất lên là $2,6\text{m}$ đi qua cổng hình parabol có phương trình $y = ax^2$. Biết khoảng cách giữa hai chân cổng là $6\text{m}$ và chiều cao tối đa của cổng là $4,5\text{m}$. Hỏi xe tải có đi qua được cổng không? Giải thích?
		\item Cho 3 số thực dương $x; y; z$ thỏa mãn $\frac{1}{x^2} + \frac{1}{y^2} + \frac{1}{z^2} = 1$. Tìm giá trị nhỏ nhất của biểu thức:
		$P = \frac{y^2z^2}{x(y^2+z^2)} + \frac{z^2x^2}{y(z^2+x^2)} + \frac{x^2y^2}{z(x^2+y^2)}$.
	\end{enumerate}
	\loigiai{}
\end{ex}

\setcounterref{ex}{0}
